\chapter{Capstone 故障模式与升级处理手册}

\section{案例导读：好的课程系统要知道什么时候该升级处理}

Capstone 中一定会出现故障：数据下载失败、学生项目越界、notebook 无法复现、AI 使用边界不清、评分争议升高。真正危险的不是出错本身，而是每个人都觉得问题重要，却没人知道谁该接手。

本章把故障模式写成 escalation playbook。它训练课程团队把问题从“抱怨”转成“分级、路由、责任人和关闭条件”，让复杂情况不会在期末集中爆炸。

\section{案例目标}

第 53 章把学期重启前的预检、权限、owner 和启动 blocker 收束成 reboot checklist。本章继续处理课程真正运行时最棘手的一层：\textbf{当 notebook 崩了、政策边界不清、评分争议升级、私有数据访问异常时，课程团队怎样判断哪些问题可以本地修、哪些必须升级、哪些应该立刻暂停 rollout。}

本章的目标有四点：

\begin{enumerate}
    \item 理解 failure mode 不是坏运气，而是课程运行里必须预先设计的分流节点；
    \item 学会检查 signal、severity、reproducibility、owner、escalation boundary 和 urgency；
    \item 学会区分“本地补丁即可”“补复现证据”“先指定事件 owner”“暂停 rollout 再修”和“升级到人类复核”五类出口；
    \item 学会把课程运行故障从个人慌乱反应变成可执行的 escalation playbook。
\end{enumerate}

\section{最危险的不是出错，而是不知道谁该接手}

一门 capstone 课程真正运行起来之后，最难的时刻往往不是学术内容，而是事故边缘时刻：学生说 notebook 在实验室电脑上跑不动，助教对评分边界理解不一致，课程团队发现某些示例可能碰到隐私或 AI-use policy 的灰区，或者 kickoff 当天一批链接突然失效。

如果团队没有 failure-mode playbook，最常见的结果不是问题本身更严重，而是大家不知道谁该先处理、什么时候该暂停、什么时候必须升级给课程负责人或政策审核人。

所以，本章的核心立场是：\textbf{课程故障管理的重点不是“永不出错”，而是让团队知道什么能本地修，什么必须升级，什么应该立刻停下来。}

\section{一个极小的 failure-mode 数据集}

配套 notebook 使用 \nolinkurl{capstone_failure_mode_escalation_demo.csv}。这是一组完全本地、教学用途的\textbf{课程故障卡片}。每一行代表一个运行中的 failure mode，包含：

\begin{itemize}
    \item \texttt{incident\_id}；
    \item \texttt{failure\_area}；
    \item \texttt{failure\_signal\_status}；
    \item \texttt{severity\_status}；
    \item \texttt{reproducibility\_status}；
    \item \texttt{owner\_status}；
    \item \texttt{escalation\_boundary\_status}；
    \item \texttt{urgency\_score}；
    \item \texttt{reference\_route}。
\end{itemize}

这里的 route 分成五类：

\begin{itemize}
    \item \texttt{local\_patch\_ready}：问题有界、可复现、owner 明确，可以本地补丁修复；
    \item \texttt{collect\_repro\_before\_escalation}：还需要复现步骤、日志或截图；
    \item \texttt{assign\_incident\_owner}：没有明确负责人，不能继续靠群体模糊响应；
    \item \texttt{pause\_rollout\_and\_patch}：问题已经扩散，应该先暂停 rollout；
    \item \texttt{escalate\_to\_human\_review}：触及政策、隐私、成绩或关键边界，必须升级。
\end{itemize}

\section{先看一个危险 baseline：只按紧急程度排序}

课程故障发生时，一个非常常见的 baseline 是只按 urgency score 排序：越紧急的先处理。

\begin{examplebox}[只按 urgency 选择当前要处理的事件]
\begin{lstlisting}[style=pythonstyle]
top4 = sorted(
    rows,
    key=lambda row: row["urgency_score"],
    reverse=True,
)[:4]
\end{lstlisting}
\end{examplebox}

这个 baseline 的问题在于，紧急不等于可直接修复。高紧急度事件里，可能混入政策问题、隐私边界问题、需要暂停 rollout 的系统性问题，或者根本还没复现成功的问题。

在本章教学数据上，只按紧急程度取前四项，真正适合本地补丁处理的精度只有 0.25，非本地修复项混入率是 0.75。表~\ref{tab:ch54_urgency_vs_escalation} 展示了结构化 escalation workflow 如何先看 boundary，再看 signal 是否已经扩散，再检查 owner 与 reproducibility：

\begin{table}[htbp]
\centering
\small
\setlength{\tabcolsep}{4pt}
\begin{tabular}{@{}p{0.28\textwidth}>{\centering\arraybackslash}p{0.18\textwidth}>{\centering\arraybackslash}p{0.18\textwidth}p{0.28\textwidth}@{}}
\toprule
方法 & 本地修复精度 & 非本地项混入率 & 教学解读 \\
\midrule
只按紧急程度取前四项 & 0.25 & 0.75 & 紧急度高的事件不一定能直接修，可能要暂停或升级 \\
结构化 escalation workflow & 1.00 & 0.00 & 先判边界和扩散，再决定本地修、暂停还是升级 \\
\bottomrule
\end{tabular}
\caption{本章 notebook 中两个最小故障处理流程的对照。课程故障管理不是 urgency 排行榜。}
\label{tab:ch54_urgency_vs_escalation}
\end{table}

\section{一个可执行的 escalation workflow}

本章 notebook 的 routing 逻辑可以写成：

\begin{examplebox}[一个最小故障升级 routing 逻辑]
\begin{lstlisting}[style=pythonstyle]
if escalation_boundary_status in {"policy", "privacy", "grading"}:
    escalate_to_human_review
elif failure_signal_status == "widespread":
    pause_rollout_and_patch
elif owner_status != "assigned":
    assign_incident_owner
elif reproducibility_status != "reproduced":
    collect_repro_before_escalation
else:
    local_patch_ready
\end{lstlisting}
\end{examplebox}

这个顺序体现了一个运行原则：\textbf{边界先于效率，扩散先于补丁，责任人先于讨论，复现先于结论。}

\section{故障升级手册应该包含什么}

一份可执行的 escalation playbook 至少应该包含：

\begin{enumerate}
    \item 事件分类：技术故障、评分争议、政策边界、隐私边界、发布错误；
    \item 升级边界：什么情况必须找课程负责人、助教主管、政策审核人；
    \item 暂停条件：什么情况要立刻暂停 rollout、提交入口或公开页面；
    \item 证据要求：日志、截图、复现步骤、受影响人数和当前 workaround；
    \item owner 规则：谁接 first response、谁修、谁复核、谁对外说明；
    \item 事后记录：事故摘要、修复方式、遗留 caveat 和下次如何避免。
\end{enumerate}

这份手册不需要长篇制度化文字，但必须让团队在出事时不用即兴发明流程。

\section{配套 notebook 做了什么}

本章 notebook 采用的是一个 teaching-first 的最小设计：

\begin{itemize}
    \item 先读取 incident table，并统计 severity、signal、owner、boundary 和 reference route；
    \item 用 urgency score 做一个 naive incident baseline；
    \item 再实现一个带 boundary、widespread signal、owner 和 reproducibility 的 escalation workflow；
    \item 输出 local-patch、collect-repro、assign-owner、pause-rollout 和 human-review 五个队列；
    \item 为每个非 local 事件生成处理前 checklist；
    \item 最后生成 escalation playbook outline 和 incident assistant prompt。
\end{itemize}

\section{AI 助手如何帮助你}

在这个案例里，AI 助手最适合帮助你：

\begin{itemize}
    \item 把散乱的学生报错、助教反馈和日志整理成 incident table；
    \item 从对话中抽取可复现信息、受影响范围和当前 workaround；
    \item 标出可能触及政策、隐私或成绩边界的事件；
    \item 为课程负责人生成一页 incident briefing；
    \item 把事后复盘写成下轮课程的 failure-mode note。
\end{itemize}

但 AI 助手不能替团队判断政策边界，也不能替代真正的事故负责人与最终决策。它能帮你更快看清局面，却不能替你承担升级责任。

\section{练习}

\begin{exercisebox}[基础练习]
把一个 \texttt{local\_patch\_ready} 事件的 \texttt{reproducibility\_status} 改成 \texttt{missing}。重新运行 notebook，观察它为什么必须先离开本地补丁队列。
\end{exercisebox}

\begin{exercisebox}[进阶练习]
新增一个 urgency 很高、但 \texttt{escalation\_boundary\_status = privacy} 的事件。预测它会落到哪个 route，并说明为什么高紧急度不能绕过人类复核。
\end{exercisebox}

\begin{exercisebox}[开放练习]
为你自己的课程项目写一页 escalation playbook。至少包含：事件分类、暂停条件、升级边界、证据要求、owner 规则和事后复盘格式。
\end{exercisebox}

\section{本章小结}

课程运行中的故障无法完全避免，但可以被更好地接住。本章把 signal、severity、reproducibility、owner、boundary 和 urgency 放进同一张 incident card。

当课程团队能把事件分成本地补丁、补复现证据、先指定事件 owner、暂停 rollout 再修和升级到人类复核，Capstone 就从“出事时靠经验顶住”变成“出事时有章可循”的课程系统。
